\section{Appendix: Lean Interface and Formalization Scaffold}

\subsection{Overview}

This appendix defines the interface between the TeX mathematical spine
and the Lean formal verification environment. The purpose is not to
provide full proofs, but to specify the types, structures, and proof
obligations required for a complete formalization of the Redux pipeline.

Each operator introduced in previous sections is represented as a Lean
structure with associated properties. This appendix serves as the
canonical reference for the Lean modules.

\subsection{Manifold Structure}

We define the manifold type as
\[
M : \mathrm{Type},
\]
equipped with a metric tensor
\[
g : M \to \mathrm{Matrix}\; 80\; 80,
\]
and a curvature tensor
\[
R : M \to \mathrm{Tensor}\; 3.
\]

Lean representation:
\begin{verbatim}
structure Manifold :=
  (dim : ℕ)
  (metric : Point → Matrix dim dim)
  (curvature : Point → Tensor 3 dim)
\end{verbatim}

\subsection{Spectral Operator}

The spectral operator is represented as a structure containing the
eigenbasis and projection map.

\[
U(x) : \mathrm{Matrix}\; 80\; 80,
\quad
s = U(x)^\top v.
\]

Lean signature:
\begin{verbatim}
structure Spectral :=
  (U : Point → Matrix 80 80)
  (project : Point → Vector 80 → Vector 80)
\end{verbatim}

\subsection{Curvature Envelope}

The curvature envelope is represented as a function
\[
\kappa : M \to \mathrm{Vector}\; 80.
\]

Lean signature:
\begin{verbatim}
structure CurvatureEnvelope :=
  (kappa : Point → Vector 80)
  (modulate : Point → Vector 80 → Vector 80)
\end{verbatim}

\subsection{Holonomy Transport}

Holonomy transport is represented as an orthogonal matrix
\[
H(x) \in \mathrm{O}(80).
\]

Lean signature:
\begin{verbatim}
structure Holonomy :=
  (H : Point → Matrix 80 80)
  (transport : Point → Vector 80 → Vector 80)
\end{verbatim}

\subsection{Topology Modulation}

The topology signature is represented as
\[
\tau : M \to \mathrm{Vector}\; 80.
\]

Lean signature:
\begin{verbatim}
structure Topology :=
  (tau : Point → Vector 80)
  (modulate : Point → Vector 80 → Vector 80)
\end{verbatim}

\subsection{Embedding Pipeline}

The full embedding pipeline is represented as a composition of the
operators defined above.

Lean signature:
\begin{verbatim}
structure Pipeline :=
  (Φ₁ : Structural)
  (Φ₂ : Spectral)
  (Φ₃ : CurvatureEnvelope)
  (Φ₄ : Holonomy)
  (Φ₅ : Topology)
  (embed : Graph → Vector 80)
\end{verbatim}

\subsection{Invariant Layer}

Each invariant is represented as a function
\[
I_k : \mathrm{Vector}\; 80 \to \mathbb{R}.
\]

Lean signature:
\begin{verbatim}
structure Invariant :=
  (eval : Vector 80 → ℝ)
\end{verbatim}

\subsection{Proof Obligations}

The following properties must be proven in Lean:

\begin{itemize}
  \item \textbf{Stability:}
    \[
    \| E(G) - E(G') \| \leq C \| G - G' \|.
    \]
  \item \textbf{Reversibility:}
    \[
    Φ_k^{-1}(Φ_k(v)) = v.
    \]
  \item \textbf{Smoothness:}
    All operators depend smoothly on the metric tensor.
  \item \textbf{Altitude monotonicity:}
    No operator lowers altitude.
\end{itemize}

These obligations define the formal verification roadmap.

\subsection{Conclusion}

This appendix provides the Lean interface required for formalizing the
Redux pipeline. It serves as the bridge between the TeX mathematical
spine and the proof environment.
